% ======================================================================
\chapter{Introdução}
\label{cap:introducao}
% ======================================================================

O GitHub é hoje a principal plataforma de desenvolvimento colaborativo. Ele
guarda o código, registra quem mudou o quê, organiza a revisão entre colegas e
ainda publica sites e roda automações. Tudo isso em volta de uma ferramenta que
vem antes dele: o \textbf{Git}.

Este guia acompanha a capacitação de Git e GitHub da liga acadêmica. A
apresentação dura duas horas; aqui está tudo o que foi dito com mais calma, os
passos completos para Windows, Linux e macOS, e os exercícios para praticar
depois.

\section{Objetivos de aprendizagem}

Ao final, o participante é capaz de:

\begin{enumerate}
  \item criar e proteger a conta no GitHub, com autenticação em dois fatores, e
        pedir os benefícios de estudante;
  \item instalar e configurar o Git e o GitHub CLI no próprio sistema;
  \item registrar alterações com commits bem escritos e trabalhar com branches;
  \item integrar trabalho com merge e resolver um conflito;
  \item contribuir com um projeto de outra pessoa por fork e Pull Request;
  \item publicar um site com o GitHub Pages e automatizar uma tarefa com o
        GitHub Actions;
  \item versionar arquivos grandes com o Git LFS e assinar commits.
\end{enumerate}

\section{O roteiro das duas horas}

\begin{table}[H]
  \centering
  \caption{Roteiro de tempo da capacitação}
  \label{tab:roteiro}
  \begin{tabularx}{\textwidth}{@{}c X c c@{}}
    \toprule
    \textbf{Bloco} & \textbf{Conteúdo} & \textbf{Início} & \textbf{Duração} \\
    \midrule
    0 & Abertura, objetivos e combinados & 00:00 & 5 min \\
    1 & A conta no GitHub, 2FA e benefícios de estudante & 00:05 & 15 min \\
    2 & Instalação, configuração e autenticação & 00:20 & 20 min \\
    3 & Git no dia a dia: commits, branches, merge e conflito & 00:40 & 20 min \\
    — & \textit{Pausa} & 01:00 & 5 min \\
    4 & Colaboração: fork, Pull Request e padrões de commit & 01:05 & 20 min \\
    5 & GitHub Pages e GitHub Actions & 01:25 & 15 min \\
    6 & Git LFS e assinatura de commits & 01:40 & 15 min \\
    7 & Encerramento e próximos passos & 01:55 & 5 min \\
    \midrule
      & \textbf{Total} & & \textbf{120 min} \\
    \bottomrule
  \end{tabularx}
\end{table}

\section{Como ler este guia}

Os comandos aparecem em blocos como este. Linhas que começam com \texttt{\#}
são comentários e não precisam ser digitadas:

\begin{bashcode}
# confere a versão do Git instalada
git --version
\end{bashcode}

Ao longo do texto, três tipos de caixa chamam a atenção:

\begin{conceito}
Uma definição importante, para guardar.
\end{conceito}

\begin{dica}
Um atalho ou uma boa prática.
\end{dica}

\begin{aviso}
Um erro comum, ou um cuidado que evita dor de cabeça.
\end{aviso}

Quando um comando traz algo entre \texttt{<sinais>}, como
\texttt{<seu-usuario>}, troque tudo, inclusive os sinais, pelo valor real.

\begin{dica}[Qual terminal usar]
No \textbf{Windows}, use o \textbf{Git Bash}, que vem com o Git (capítulo
\ref{cap:ambiente}). Os comandos deste guia, como \texttt{ls}, \texttt{touch}
e \texttt{mkdir -p}, são do Bash e não funcionam no Prompt de Comando. No
\textbf{Linux} e no \textbf{macOS}, use o Terminal do sistema. A única exceção
é a instalação pelo \texttt{winget}, feita no PowerShell.
\end{dica}

\section{Pré-requisitos}

\begin{itemize}
  \item Um computador com Windows 10 (versão 1809 ou mais nova), Windows 11,
        Linux ou macOS, com permissão para instalar programas.
  \item Um endereço de e-mail. Para os benefícios de estudante, é melhor ter o
        e-mail acadêmico.
  \item Um celular com um aplicativo autenticador (Google Authenticator,
        Microsoft Authenticator ou similar), para a autenticação em dois
        fatores.
\end{itemize}

Não é preciso saber programar. Os exercícios usam arquivos prontos, que estão
na pasta \arq{materials/} do repositório.
